\documentclass[11pt]{article}

\usepackage[letterpaper,margin=1in]{geometry}
\usepackage{amsmath,amssymb,amsthm,mathtools,bm}
\usepackage{booktabs}
\usepackage{enumitem}
\usepackage[colorlinks=true,linkcolor=blue,urlcolor=blue]{hyperref}
\usepackage{microtype}

\numberwithin{equation}{section}

\newtheorem{theorem}{Theorem}[section]
\newtheorem{lemma}[theorem]{Lemma}
\newtheorem{proposition}[theorem]{Proposition}
\newtheorem{corollary}[theorem]{Corollary}
\newtheorem{remark}[theorem]{Remark}

\DeclareMathOperator{\Tr}{Tr}
\DeclareMathOperator{\diag}{diag}
\DeclareMathOperator{\Lip}{Lip}
\newcommand{\R}{\mathbb R}
\newcommand{\E}{\mathbb E}
\newcommand{\Id}{\mathrm{Id}}
\newcommand{\op}{\mathrm{op}}
\newcommand{\loc}{\mathrm{loc}}

\title{A Bounded-Hessian-Lipschitz Counterexample to the\\
Dimension-Free Logarithmic Local Rate}
\author{}
\date{July 14, 2026}

\begin{document}

\maketitle

\begin{abstract}
We give an option-(B) counterexample to the Hessian-Lipschitz local-rate
conjecture.  A two-scale radial-softplus construction produces smooth genuine
Hessians, already in equilibrium-whitened coordinates, for which
\[
 \kappa\asymp \frac d{\sigma^2},\qquad
 \gamma_{\loc}\asymp\frac{\sigma}{\sqrt d}\asymp\kappa^{-1/2},\qquad
 L_H\asymp\frac{\sqrt d}{\sigma^2},\qquad
 \tau_H^2\gtrsim\frac{\sqrt d}{\sigma}\asymp\sqrt\kappa.
\]
Taking $(d,\sigma)=(n^4,n)$ gives $L_H\le\Lambda_0$ for one absolute
constant, $\kappa\asymp n^2$, and
$\gamma_{\loc}\asymp n^{-1}\asymp\kappa^{-1/2}$.  Thus bounded whitened
Hessian-Lipschitz modulus does not rescue a logarithmic local rate.  The proof
uses exact Gaussian isotropization and a literally Fisher--Rao-normalized
affine Rayleigh witness.  No Frobenius/operator conversion, symmetry-only
surrogate, or noncommutative Loewner transfer occurs.
\end{abstract}

\section{Main result and norm convention}

For a $C^3$ potential define its Hessian operator-Lipschitz modulus by
\[
 L_H(V):=\sup_{x\ne y}
 \frac{\|\nabla^2V(x)-\nabla^2V(y)\|_{\op}}{\|x-y\|}.
\]
Since the third derivative is symmetric,
\begin{equation}\label{eq:tensor-lip}
 L_H(V)=\sup_x\sup_{\|h\|=\|u\|=1}
 |\nabla^3V(x)[h,u,u]|.
\end{equation}
Indeed, the right side is
$\sup_{x,\|h\|=1}\|D(\nabla^2V)(x)[h]\|_{\op}$; the fundamental theorem of
calculus proves one inequality and directional difference quotients prove the
other.  Thus every $L_H$ below is the tensor/operator quantity in (H2), not a
Frobenius norm.
For a symmetric trilinear form $\mathcal T$, we also write
$\|\mathcal T\|_{\rm inj}:=\sup_{\|h\|=\|u\|=\|v\|=1}
|\mathcal T[h,u,v]|$.

\begin{theorem}[Bounded-$L_H$ square-root counterexample]\label{thm:main}
There are numerical constants $c,C>0$ and $n_0\in\mathbb N$ such that for
every integer $n\ge n_0$ there is a $C^\infty$, strongly convex potential
$V_n:\R^{n^4+1}\to\R$ satisfying, for
$Z\sim\mathcal N(0,I_{n^4+1})$,
\begin{equation}\label{eq:equilibrium-main}
 \E\nabla V_n(Z)=0,\qquad \E\nabla^2V_n(Z)=I.
\end{equation}
If $\alpha_n,\beta_n$ are its optimal scalar curvature constants and
$\kappa_n=\beta_n/\alpha_n$, then
\begin{equation}\label{eq:main-scales}
 \frac c n\le\alpha_n\le\frac Cn,\qquad
 cn\le\beta_n\le Cn,\qquad
 cn^2\le\kappa_n\le Cn^2.
\end{equation}
Its whitened Hessian-Lipschitz modulus and linearized gap obey
\begin{equation}\label{eq:main-L-gap}
 c\le L_H(V_n)\le C,
 \qquad
 \frac c n\le\gamma_{\loc}(V_n)\le\frac Cn,
\end{equation}
and its first-Hermite Hessian coupling satisfies
\begin{equation}\label{eq:main-tau}
 \tau_H(V_n)^2\ge cn.
\end{equation}
Consequently
\[
 \gamma_{\loc}(V_n)\asymp\kappa_n^{-1/2},
 \qquad
 \gamma_{\loc}(V_n)^{-1}\asymp\sqrt{\kappa_n}\gg\log\kappa_n.
\]
In particular, for an absolute $\Lambda_0$ the whole sequence belongs to
$\mathcal C(\kappa_n,\Lambda_0)$, while
$\gamma_{\loc}(V_n)(1+\log\kappa_n)\to0$.
\end{theorem}

The construction and proof occupy Sections~\ref{sec:family}--\ref{sec:tau}.

\section{The two-scale radial-softplus family}\label{sec:family}

Let $d$ and $\sigma$ satisfy
\begin{equation}\label{eq:parameter-range}
 \sigma\ge1,\qquad d\ge C_0\sigma^2,
\end{equation}
where $C_0$ is a sufficiently large numerical constant.  Write
$z=(t,y)\in\R\times\R^d$ and set
\begin{equation}\label{eq:scaled-profiles}
 r(y):=\sqrt{\sigma^2+\|y\|^2},\qquad
 R:=\sqrt{\sigma^2+d},\qquad
 F_\sigma(s):=\sigma\log(1+e^{s/\sigma}).
\end{equation}
Put
\[
 p(s):=F_\sigma'(s),\qquad
 q(s):=F_\sigma''(s)=\frac1\sigma p(s)(1-p(s)).
\]
Thus $0<p<1$, $0<q\le(4\sigma)^{-1}$, and
$|F_\sigma'''|\le(4\sigma^2)^{-1}$.

For $0\le\delta\le1$, define
\begin{equation}\label{eq:Udelta}
 s_\delta(t,y):=r(y)-R-\delta t,\qquad
 U_\delta(t,y):=F_\sigma(s_\delta(t,y)),\qquad
 A_\delta:=\nabla^2U_\delta.
\end{equation}
Writing $S=\|y\|^2$,
\begin{equation}\label{eq:radial-derivatives}
 v_\delta:=\nabla s_\delta=(-\delta,y/r),\qquad
 D(y):=\nabla_y^2r=\frac{I_d}{r}-\frac{yy^\top}{r^3}.
\end{equation}
The exact Hessian is
\begin{equation}\label{eq:exact-Hessian}
 A_\delta=q(s_\delta)v_\delta v_\delta^\top
 +p(s_\delta)\begin{pmatrix}0&0\\0&D(y)\end{pmatrix}.
\end{equation}
Both terms are positive semidefinite.  The tangential eigenvalue of $D$ is
$r^{-1}$ and its radial eigenvalue is $\sigma^2r^{-3}$.  Since
$\|v_\delta\|^2\le2$,
\begin{equation}\label{eq:A-bound}
 0\preceq A_\delta\preceq\frac{3}{2\sigma}I_{d+1}.
\end{equation}

\section{Exact Gaussian isotropization}\label{sec:isotropization}

Let $T\sim\mathcal N(0,1)$ and $Y\sim\mathcal N(0,I_d)$ be independent.
In the following expectations $s_\delta=s_\delta(T,Y)$ and $r=r(Y)$.
Define
\begin{align}
 Q_{d,\sigma}(\delta)&:=\E q(s_\delta),\label{eq:Qdef}\\
 M_{d,\sigma}(\delta)&:=\frac1d\E\left[
 q(s_\delta)\frac S{r^2}
 +p(s_\delta)\left(\frac{d-1}{r}+\frac{\sigma^2}{r^3}\right)
 \right].\label{eq:Mdef}
\end{align}
Rotation invariance and oddness of the mixed block give
\begin{equation}\label{eq:mean-A}
 \E A_\delta
 =\diag(\delta^2Q_{d,\sigma}(\delta),M_{d,\sigma}(\delta)I_d).
\end{equation}

\begin{lemma}[Uniform scalar estimates and the isotropizing root]
\label{lem:isotropization}
Under~\eqref{eq:parameter-range}, uniformly for $0\le\delta\le1$,
\begin{equation}\label{eq:QM-bounds}
 \frac c\sigma\le Q_{d,\sigma}(\delta)\le\frac1{4\sigma},
 \qquad
 \frac cR\le M_{d,\sigma}(\delta)\le\frac CR.
\end{equation}
If $C_0$ is large enough, the continuous equation
\begin{equation}\label{eq:root-equation}
 \delta^2Q_{d,\sigma}(\delta)=M_{d,\sigma}(\delta)
\end{equation}
has a root $\delta_{d,\sigma}\in(0,1)$.  For any such root, with
$m_{d,\sigma}$ denoting the common value,
\begin{equation}\label{eq:root-scales}
 \delta_{d,\sigma}^2\asymp\frac\sigma R,\qquad
 m_{d,\sigma}\asymp\frac1R,\qquad
 \E A_{\delta_{d,\sigma}}=m_{d,\sigma}I_{d+1}.
\end{equation}
\end{lemma}

\begin{proof}
On the event
\[
 \mathcal E_d:=\{|S-d|\le2\sqrt d,\ |T|\le1\},
\]
Chebyshev's inequality and independence give
$\mathbb P(\mathcal E_d)\ge c_0>0$.  Moreover
\[
 |r-R|=\frac{|S-d|}{r+R}\le2,\qquad |s_\delta|\le3.
\]
Since $\sigma\ge1$, both $p(s_\delta)$ and
$\sigma q(s_\delta)$ are bounded below on this event.  This proves the
lower bound for $Q$.  On the same event,
\[
 \Tr D=\frac{d-1}{r}+\frac{\sigma^2}{r^3}\ge c\frac dR,
\]
which proves the lower bound for $M$.

The upper bound for $Q$ is immediate.  For $d>2$,
\begin{equation}\label{eq:inverse-radius}
 \E\frac1r\le\E S^{-1/2}
 \le(\E S^{-1})^{1/2}=\frac1{\sqrt{d-2}}\le\frac CR.
\end{equation}
The first summand in~\eqref{eq:Mdef} is at most $(4\sigma d)^{-1}$,
and the second is at most $(1+d^{-1})\E r^{-1}$.  This proves
\eqref{eq:QM-bounds}.

Continuity follows by dominated convergence.  The left side minus the right
side of~\eqref{eq:root-equation} is negative at zero and positive at one
once $R/\sigma$ is a sufficiently large numerical constant.  Hence a root
exists.  At any root,~\eqref{eq:QM-bounds} gives
\eqref{eq:root-scales}, and~\eqref{eq:mean-A} gives exact isotropy.
\end{proof}

Fix a root, abbreviate $U=U_{\delta_{d,\sigma}}$, $A=\nabla^2U$, and set
\begin{equation}\label{eq:alpha-c}
 \alpha:=\frac\sigma R,\qquad
 c_{d,\sigma}:=\frac{1-\alpha}{m_{d,\sigma}}\asymp R.
\end{equation}
Define
\begin{equation}\label{eq:Vdef}
 V_{d,\sigma}(z)
 :=\frac\alpha2\|z\|^2+c_{d,\sigma}U(z)+b_{d,\sigma}^\top z,
 \qquad
 b_{d,\sigma}:=-c_{d,\sigma}\E\nabla U(Z).
\end{equation}
The gradient of $U$ is bounded, so $b_{d,\sigma}$ is finite.  It is a
multiple of the $t$ direction.  The Hessian
\begin{equation}\label{eq:Wdef}
 W_{d,\sigma}=\nabla^2V_{d,\sigma}=\alpha I+c_{d,\sigma}A
\end{equation}
satisfies, exactly,
\begin{equation}\label{eq:exact-equilibrium}
 \E\nabla V_{d,\sigma}(Z)=0,\qquad \E W_{d,\sigma}(Z)=I.
\end{equation}

The optimal lower curvature is $\alpha$, since $A(t,y)\to0$ as
$t\to+\infty$ for fixed $y$.  Equation~\eqref{eq:A-bound} gives the upper
curvature bound.  Conversely, at $y=0$ and with $t$ chosen so that
$s_\delta=0$, the transverse block of $A$ is $I_d/(2\sigma)$.  Thus the
optimal upper curvature $\beta$ satisfies
\begin{equation}\label{eq:beta-kappa}
 \alpha+\frac{c_{d,\sigma}}{2\sigma}\le\beta
 \le\alpha+\frac{3c_{d,\sigma}}{2\sigma},\qquad
 \beta\asymp\frac R\sigma,\qquad
 \kappa:=\frac\beta\alpha\asymp\frac{R^2}{\sigma^2}\asymp\frac d{\sigma^2}.
\end{equation}
The potential is $C^\infty$ and $\alpha$-strongly convex, hence normalizable.
It is a genuine Hessian and its complete compatibility hierarchy holds
classically.

\section{The exact Hessian-Lipschitz scale}\label{sec:Lip}

The third derivative of $r$ is
\begin{align}\label{eq:D3r}
 \nabla^3r[h,u,v]
={}&-\frac{(y\cdot h)(u\cdot v)+(y\cdot u)(h\cdot v)
 +(y\cdot v)(h\cdot u)}{r^3}\\
 &+3\frac{(y\cdot h)(y\cdot u)(y\cdot v)}{r^5}.
\end{align}
Consequently
\begin{equation}\label{eq:s-derivative-bounds}
 \|\nabla s_\delta\|\le\sqrt2,\qquad
 \|\nabla^2s_\delta\|_{\op}\le\sigma^{-1},\qquad
 \|\nabla^3s_\delta\|_{\rm inj}\le6\sigma^{-2}.
\end{equation}
The exact chain rule is
\begin{align}\label{eq:D3U}
 \nabla^3U[h,u,v]={}&F_\sigma'''(s)\,ds[h]ds[u]ds[v]\\
 &+q(s)\{\nabla^2s[h,u]ds[v]+\nabla^2s[h,v]ds[u]
 +\nabla^2s[u,v]ds[h]\}\\
 &+p(s)\nabla^3s[h,u,v].
\end{align}
Equations~\eqref{eq:s-derivative-bounds}--\eqref{eq:D3U} give
\begin{equation}\label{eq:D3U-upper}
 \sup_z\|\nabla^3U(z)\|_{\rm inj}
 \le\frac{6+5\sqrt2/4}{\sigma^2}<\frac8{\sigma^2}.
\end{equation}
For the reverse bound, fix $y=(\sigma/2)e_1$ and send $t\to-\infty$.
Then $p(s)\to1$ and $q(s),F_\sigma'''(s)\to0$, so
$\nabla^3U\to\nabla^3r$.  Directly from~\eqref{eq:D3r},
\[
 |\nabla^3r[e_1,e_1,e_1]|=\frac{c_1}{\sigma^2}
\]
for a numerical $c_1>0$.  Since the quadratic and affine terms in
\eqref{eq:Vdef} have zero third derivative,
\begin{equation}\label{eq:LH-frontier}
 L_H(V_{d,\sigma})\asymp\frac{c_{d,\sigma}}{\sigma^2}
 \asymp\frac R{\sigma^2}\asymp\frac{\sqrt d}{\sigma^2}.
\end{equation}
This is a direct injective-tensor estimate; no dimension-dependent
Frobenius conversion is present.

\section{Compatible Rayleigh identity and the inverse gap}\label{sec:gap}

If $X$ is the physical covariance perturbation, its packed coordinate is
$X/\sqrt2$ and its Fisher norm is
\begin{equation}\label{eq:Fisher-norm}
 \|(u,X)\|_\star^2:=\|u\|^2+\frac12\|X\|_F^2.
\end{equation}

\begin{lemma}[Compatible Bochner--Rayleigh identity]\label{lem:rayleigh}
Let $V\in C^4$ and suppose that its Hessian $W=\nabla^2V$ and the first two
derivatives of $W$ have at most polynomial growth.  If $\E W=I$, then the
packed generator in the question satisfies
\begin{equation}\label{eq:rayleigh-identity}
 4\mathcal Q_W(u,X)
 =\E[(2u+XZ)^\top W(Z)(2u+XZ)]+\|X\|_F^2.
\end{equation}
If $W\succeq\alpha I$ and $0<\alpha\le1$, then
\begin{equation}\label{eq:alpha-floor}
 \mathcal Q_W(u,X)\ge\alpha\|(u,X)\|_\star^2,\qquad
 \gamma_{\loc}(W)\ge\alpha.
\end{equation}
\end{lemma}

\begin{proof}
Expanding the displayed block operator in the packed coordinate gives
\[
 \mathcal Q_W(u,X)=\|u\|^2+u^\top T[X]
 +\frac12\|X\|_F^2+\frac14\langle X,H[X]\rangle_F.
\]
Gaussian integration by parts and third-derivative symmetry give
\[
 \E[u^\top W XZ]=u^\top T[X].
\]
In components, double Gaussian integration by parts gives
\[
 \E[Z_aZ_bW_{ij}(Z)]
 =\delta_{ab}\delta_{ij}+\E[\partial_{ab}W_{ij}(Z)].
\]
Full fourth-derivative symmetry then gives
\[
 X_{ia}X_{jb}\E[\partial_{ab}W_{ij}]
 =X_{ij}X_{ab}\E[\partial_{ab}W_{ij}],
\]
and therefore
\[
 \E[(XZ)^\top W(XZ)]
 =\|X\|_F^2+\langle X,H[X]\rangle_F.
\]
Substitution proves~\eqref{eq:rayleigh-identity}.  The displayed family has a
bounded Hessian and bounded derivatives of the Hessian of every order because
$r\ge\sigma>0$, so all steps are classical.  Finally,
\[
 4\mathcal Q_W(u,X)
 \ge4\alpha\|u\|^2+(1+\alpha)\|X\|_F^2
 \ge4\alpha\|(u,X)\|_\star^2,
\]
which proves~\eqref{eq:alpha-floor}.
\end{proof}

For the matching upper bound, put $P_d:=\diag(0,I_d)$ and first take the
unnormalized pair
\begin{equation}\label{eq:witness}
 u_0=e_t,\qquad
 X_0=\lambda P_d,\qquad
 \lambda:=\frac{2\delta_{d,\sigma}R}{d}.
\end{equation}
Its literal Fisher--Rao normalization is
\begin{equation}\label{eq:normalized-witness}
 N_F^2:=1+\frac{d\lambda^2}{2},\qquad
 \widehat u:=\frac{u_0}{N_F},\qquad
 \widehat X:=\frac{X_0}{N_F},\qquad
 \|(\widehat u,\widehat X)\|_\star=1.
\end{equation}
The covariance coordinate in the packed Euclidean display of $L_\star$ is
$\widehat X/\sqrt2$.  For the unnormalized pair,
$2u_0+X_0Z=(2,\lambda Y)$ and
\begin{equation}\label{eq:X-cost}
 \|X_0\|_F^2=d\lambda^2
 =\frac{4\delta_{d,\sigma}^2R^2}{d}\le C\frac\sigma R=C\alpha,
 \qquad N_F^2=1+O(\alpha).
\end{equation}
Let $g(S)=S/r$.  From~\eqref{eq:exact-Hessian}, exactly,
\begin{equation}\label{eq:exact-A-energy}
 (2,\lambda Y)^\top A(2,\lambda Y)
 =q(s)\lambda^2\left(g(S)-\frac dR\right)^2
 +p(s)\lambda^2\frac{\sigma^2S}{r^3}.
\end{equation}
The factorization
\begin{equation}\label{eq:residual-factorization}
 g(S)-\frac dR
 =(S-d)\frac{\sigma^2+rR}{rR(r+R)}
\end{equation}
implies
\begin{equation}\label{eq:residual-moment}
 \left|g(S)-\frac dR\right|\le\frac{|S-d|}{R},\qquad
 \E\left(g(S)-\frac dR\right)^2\le\frac{2d}{R^2}\le2.
\end{equation}
Also, by~\eqref{eq:inverse-radius},
\begin{equation}\label{eq:radial-moment}
 \E\frac S{r^3}\le\E\frac1r\le\frac C R.
\end{equation}
Using $q\le(4\sigma)^{-1}$,
$\delta_{d,\sigma}^2\asymp\sigma/R$, and
$c_{d,\sigma}\asymp R$, equations
\eqref{eq:exact-A-energy}--\eqref{eq:radial-moment} give
\begin{equation}\label{eq:scaled-A-energy}
 c_{d,\sigma}\E[(2,\lambda Y)^\top A(2,\lambda Y)]
 \le C\left(\frac1d+\frac{\sigma^3R}{d^2}\right)
 \le C\frac\sigma R=C\alpha.
\end{equation}
Moreover
\begin{equation}\label{eq:floor-energy}
 \alpha\E\|(2,\lambda Y)\|^2=\alpha(4+d\lambda^2)\le C\alpha.
\end{equation}
Substitution into~\eqref{eq:rayleigh-identity} proves
$\mathcal Q_W(u_0,X_0)\le C\alpha$.  By homogeneity and
\eqref{eq:normalized-witness},
\begin{equation}\label{eq:gap-frontier}
 \gamma_{\loc}(V_{d,\sigma})
 \le\mathcal Q_W(\widehat u,\widehat X)
 =\frac{\mathcal Q_W(u_0,X_0)}{N_F^2}
 \le C\frac\sigma R.
\end{equation}
The alpha floor gives the reverse bound.  Together with
\eqref{eq:beta-kappa},
\begin{equation}\label{eq:gap-exact-frontier}
 \gamma_{\loc}(V_{d,\sigma})\asymp\frac\sigma R
 \asymp\kappa^{-1/2}.
\end{equation}

\subsection{Symmetry-sector check}

The construction is $O(d)$-invariant in the $y$ variables.  Accordingly the
packed mean--covariance space splits orthogonally into a three-dimensional
trivial sector
\[
 \operatorname{span}\{e_t;\ e_te_t^\top;\ P_d/\sqrt d\},
\]
a two-dimensional standard sector
$\operatorname{span}\{e_j;(e_te_j^\top+e_je_t^\top)/\sqrt2\}$ repeated $d$
times, and a transverse-traceless scalar sector of multiplicity
$d(d+1)/2-1$.  The multiplicities sum to
$(d+1)+(d+1)(d+2)/2$, the dimension of the full packed space.

For an exact check, write
\[
 W=\begin{pmatrix}a&b y^\top\\b y&c_0I_d+e\,yy^\top\end{pmatrix}
\]
with
\[
 a=\alpha+c_{d,\sigma}\delta^2q,\quad
 b=-c_{d,\sigma}\delta q/r,\quad
 c_0=\alpha+c_{d,\sigma}p/r,\quad
 e=c_{d,\sigma}(q/r^2-p/r^3).
\]
Let $E=e_te_t^\top$, $P=P_d/\sqrt d$, and set
\begin{align*}
 t_E&=\E[Ta],&
 t_P&=\frac1{\sqrt d}\E[T(dc_0+eS)]
 =-\frac{c_{d,\sigma}\delta}{\sqrt d}\E[qS/r],\\
 h_{EE}&=\E[a(T^2-1)],&
 h_{EP}&=\E[a(S-d)/\sqrt d]
 =\frac1{\sqrt d}\E[(dc_0+eS)(T^2-1)],\\
 h_{PP}&=\E[(dc_0+eS)(S-d)/d].
\end{align*}
In the orthonormal packed basis $(e_t,E,P)$, the trivial block is exactly
\begin{equation}\label{eq:trivial-sector}
 L_{\rm triv}=
 \begin{pmatrix}
 1&t_E/\sqrt2&t_P/\sqrt2\\
 t_E/\sqrt2&1+h_{EE}/2&h_{EP}/2\\
 t_P/\sqrt2&h_{EP}/2&1+h_{PP}/2
 \end{pmatrix}.
\end{equation}
The normalized witness~\eqref{eq:normalized-witness} has the exact sector
coordinate
\begin{equation}\label{eq:witness-sector-coordinate}
 \frac{(1,0,\lambda\sqrt{d/2})}
 {\sqrt{1+d\lambda^2/2}}
\end{equation}
in~\eqref{eq:trivial-sector}.  Its block Rayleigh quotient is therefore
identically the full-space quotient evaluated in
\eqref{eq:exact-A-energy}--\eqref{eq:gap-frontier}; no multiplicity or
physical-versus-packed coordinate is being suppressed.

\section{First-Hermite certificate}\label{sec:tau}

Let
\begin{equation}\label{eq:Btest}
 B:=\diag(0,I_d/\sqrt d),\qquad \|B\|_F=1.
\end{equation}
Hessian compatibility and two Gaussian integrations by parts give
\begin{align}
 e_t^\top T[B]
 &=\E[t\Tr(BW_{d,\sigma})]
 =\E[(BZ)^\top W_{d,\sigma}e_t]\notag\\
 &=-\frac{c_{d,\sigma}\delta_{d,\sigma}}{\sqrt d}
 \E\left[q(s)\frac S r\right].\label{eq:T-certificate}
\end{align}
On the event $\mathcal E_d$ from Lemma~\ref{lem:isotropization},
$q(s)\ge c/\sigma$ and $S/r\ge c\sqrt d$.  Hence
\begin{equation}\label{eq:tau-frontier}
 \tau_H^2\ge|e_t^\top T[B]|^2
 \ge c\frac{c_{d,\sigma}^2\delta_{d,\sigma}^2}{\sigma^2}
 \ge c\frac R\sigma
 \asymp\sqrt\kappa.
\end{equation}
This explicitly realizes the central Frobenius trap rather than using it:
the averaged $t$-derivative of the transverse Hessian has tiny individual
eigenvalues but rank $d$, so its Frobenius norm is large even though the
pointwise injective third-tensor norm is $L_H\asymp R/\sigma^2$.

Combining~\eqref{eq:beta-kappa}, \eqref{eq:LH-frontier},
\eqref{eq:gap-exact-frontier}, and~\eqref{eq:tau-frontier} yields the complete
two-parameter frontier
\begin{equation}\label{eq:complete-frontier}
 \boxed{
 \kappa\asymp\frac d{\sigma^2},\qquad
 L_H\asymp\frac{\sqrt d}{\sigma^2},\qquad
 \gamma_{\loc}\asymp\frac\sigma{\sqrt d}\asymp\kappa^{-1/2},\qquad
 \tau_H^2\gtrsim\frac{\sqrt d}{\sigma}\asymp\sqrt\kappa.}
\end{equation}
The specialization $(d,\sigma)=(n^4,n)$ proves
Theorem~\ref{thm:main}.  The alternative $(d,\sigma)=(n^6,n^2)$ even has
$L_H\asymp n^{-1}\to0$ with the same $\kappa\asymp n^2$ and gap scale.
More generally, keeping
$d\asymp\kappa\sigma^2$ leaves the square-root gap unchanged while
$L_H\asymp\sqrt\kappa/\sigma$ can be made arbitrarily small by increasing
the dimension.

\section{An explicit option-(B) sequence}\label{sec:explicit-sequence}

We record loose numerical constants so that the class quantifiers do not rest
on asymptotic notation.  For every integer $n\ge15000$, set
\begin{equation}\label{eq:explicit-slice}
 d_n=n^4,\qquad \sigma_n=n,\qquad
 R_n=\sqrt{d_n+\sigma_n^2},\qquad N_\theta^{(n)}=d_n+1.
\end{equation}
Choose the smallest root $\delta_n\in(0,1)$ of
\eqref{eq:root-equation}, put
\[
 m_n:=\delta_n^2Q_{d_n,\sigma_n}(\delta_n)
     =M_{d_n,\sigma_n}(\delta_n),\quad
 \alpha_n:=\frac{\sigma_n}{R_n},\quad
 c_n:=\frac{1-\alpha_n}{m_n},
\]
and define $V_n$ by~\eqref{eq:Vdef}, with these specialized parameters.
Also write
\begin{equation}\label{eq:explicit-specialized-notation}
 A_n:=A_{\delta_n},\qquad W_n:=\nabla^2V_n,\qquad
 \lambda_n:=\frac{2\delta_nR_n}{d_n},\qquad
 N_{F,n}^2:=1+\frac{d_n\lambda_n^2}{2},
\end{equation}
and set
$\widehat u_n=e_t/N_{F,n}$,
$\widehat X_n=\lambda_nP_{d_n}/N_{F,n}$.
Then $\|(\widehat u_n,\widehat X_n)\|_\star=1$ exactly, and its packed
covariance coordinate is $\widehat X_n/\sqrt2$.

Here are fully explicit bounds behind the root.  On
$\mathcal E_n=\{|S-d_n|\le2\sqrt{d_n},\ |T|\le1\}$, Chebyshev's inequality,
independence, and $\mathbb P(|T|\le1)>1/3$ give
$\mathbb P(\mathcal E_n)>1/6$.  On this event $|s_\delta|\le3$.  Since
Euler's number satisfies $\mathrm e<3$, uniformly for $0\le\delta\le1$,
\begin{equation}\label{eq:explicit-QM}
 \frac1{4704\sigma_n}\le Q_{d_n,\sigma_n}(\delta)
 \le\frac1{4\sigma_n},\qquad
 \frac1{672R_n}\le M_{d_n,\sigma_n}(\delta)
 \le\frac3{R_n}.
\end{equation}
For the lower bounds use $p\ge1/28$, $q\ge1/(784\sigma_n)$, and
$\Tr D\ge d_n/(4R_n)$ on $\mathcal E_n$.  For the upper bound on $M$, use
\[
 \frac1{d_n}p\left(\frac{d_n-1}{r}
 +\frac{\sigma_n^2}{r^3}\right)\le\frac1r,\qquad
 \frac1{4\sigma_nd_n}\le\frac1{R_n},\qquad
 \E r^{-1}\le(\E S^{-1})^{1/2}
 =(d_n-2)^{-1/2}\le2R_n^{-1}.
\]
Since $R_n/\sigma_n=\sqrt{n^2+1}>14112$, the left side minus the right
side of~\eqref{eq:root-equation} is negative at zero and positive at one.
Thus the specified root exists, and~\eqref{eq:explicit-QM} gives
\begin{equation}\label{eq:explicit-root-bounds}
 \frac{\sigma_n}{168R_n}\le\delta_n^2
 \le14112\frac{\sigma_n}{R_n},\qquad
 \frac1{672R_n}\le m_n\le\frac3{R_n},\qquad
 \frac{R_n}{6}\le c_n\le672R_n.
\end{equation}
In particular, the two equalities in~\eqref{eq:exact-equilibrium} remain
exact; the inequalities in this paragraph are used only to audit scales.

Define the explicit valid upper curvature and class ceiling
\begin{equation}\label{eq:explicit-kappa-ceiling}
 \widetilde\beta_n:=1009\frac{R_n}{\sigma_n},\qquad
 \overline\kappa_n:=\frac{\widetilde\beta_n}{\alpha_n}
 =1009(n^2+1).
\end{equation}
Equations~\eqref{eq:A-bound} and~\eqref{eq:explicit-root-bounds} give
$\alpha_nI\preceq W_n\preceq\widetilde\beta_nI$.  Moreover the direct
chain-rule estimate~\eqref{eq:D3U-upper} can be read with constant eight, so
\begin{equation}\label{eq:explicit-Lambda}
 L_H(V_n)\le\frac{8c_n}{\sigma_n^2}
 \le5376\frac{R_n}{\sigma_n^2}\le8064=:\Lambda_0.
\end{equation}
Consequently $(N_\theta^{(n)},V_n)\in
\mathcal C(\overline\kappa_n,\Lambda_0)$ and
$\overline\kappa_n\to\infty$.

For completeness, the constants in the normalized Rayleigh calculation can
also be made numerical.  Equations~\eqref{eq:explicit-root-bounds},
\eqref{eq:exact-A-energy}, and~\eqref{eq:residual-moment} give
\[
 d_n\lambda_n^2\le112896\alpha_n,\qquad
 \E[(2,\lambda_nY)^\top A_n(2,\lambda_nY)]
 \le\frac{\lambda_n^2}{2\sigma_n}
 +2\frac{\lambda_n^2\sigma_n^2}{R_n}.
\]
Using $R_n^2/d_n\le2$, $R_n^2/d_n^2\le\alpha_n$, and
$\sigma_n^3R_n/d_n^2\le\alpha_n$, the two terms after multiplication by
$c_n$ are at most $18966528\alpha_n$ and $75866112\alpha_n$,
respectively.  Thus
\[
 c_n\E[(2,\lambda_nY)^\top A_n(2,\lambda_nY)]
 \le94832640\alpha_n.
\]
Together with the floor energy and the final $\|X\|_F^2$ term in
\eqref{eq:rayleigh-identity}, which are at most $112900\alpha_n$ and
$112896\alpha_n$, respectively, this yields
\[
 4\mathcal Q_{W_n}(e_t,\lambda_nP_{d_n})<200000000\alpha_n.
\]
For the literally normalized witness~\eqref{eq:normalized-witness}, therefore,
\begin{equation}\label{eq:explicit-psi-bound}
 \gamma_{\loc}(V_n)\le50000000\alpha_n
 \le\frac{1600000000}{\sqrt{\overline\kappa_n}}.
\end{equation}
If $\kappa_n^{\rm opt}$ denotes the optimal post-whitening condition number,
the point $y=0$, $s=0$ and~\eqref{eq:explicit-kappa-ceiling} give
\begin{equation}\label{eq:optimal-kappa-comparison}
 \frac{n^2+1}{12}\le\kappa_n^{\rm opt}
 \le1009(n^2+1).
\end{equation}
The same last bound in~\eqref{eq:explicit-psi-bound}, with the same numerical
constant, therefore holds with $\kappa_n^{\rm opt}$ in place of
$\overline\kappa_n$.  Hence outcome (B) holds with
\begin{equation}\label{eq:psi}
 \psi(\kappa):=1600000000\,\kappa^{-1/2},\qquad
 \psi(\kappa)(1+\log\kappa)\longrightarrow0.
\end{equation}

\section{Final guardrail audit}

The construction requires no coordinate transfer: it is defined with
$Z\sim N(0,I)$ and satisfies $\E W(Z)=I$ exactly.  Thus every
$\alpha,\beta,\kappa,L_H$ above is already post-whitening.  The remaining
guardrails are met as follows.
\begin{enumerate}[leftmargin=2em,label=\arabic*.,itemsep=.2em,topsep=.3em]
 \item Equation~\eqref{eq:exact-Hessian} is the pointwise Hessian of the
 displayed $C^\infty$ potential.  Full third- and fourth-order compatibility,
 not first-Hermite symmetry alone, is used in Lemma~\ref{lem:rayleigh} and
 \eqref{eq:T-certificate}.
 \item No commutation-free Loewner transfer occurs.  The inverse gap is the
 scalar Rayleigh calculation
 \eqref{eq:exact-A-energy}--\eqref{eq:gap-exact-frontier}.
 \item Bounds and isotropic mean are not treated as sufficient abstract data.
 Exact isotropy comes from~\eqref{eq:root-equation}, and the small energy comes
 from the compatible radial and mixed blocks of~\eqref{eq:exact-Hessian}.
 \item Both $\kappa$ and $L_H$ are post-whitening quantities; no raw-coordinate
 estimate is imported.
 \item The Hessian and every derivative of the Hessian are bounded, so every
 Gaussian integration by parts used here is classical.  No unproved weak
 extension is needed.
 \item The inverse gap is $\Theta(\sqrt\kappa)$, genuinely larger than
 $O(\log\kappa)$, even when $L_H\to0$.
 \item No statement about a nonlinear attraction radius is made.
 \item At no point is an operator bound on $\nabla^3V$ converted into a
 Frobenius bound.  The large Frobenius coupling is instead certified directly
 by the normalized test~\eqref{eq:Btest}.
\end{enumerate}

\end{document}
